\chapter{Trajectory pre-processing}
\label{ch:preproc}

The pre-processing turns each IGC file into a trajectory fit for analysis. This
chapter covers its first two stages: reading the fixes and choosing the altitude
channel.

\section{Reading the fixes}
\label{sec:parsing}

Each \texttt{B} record gives a fix $(t_i,\varphi_i,\lambda_i)$, with latitude and
longitude in degrees, together with the validity flag and the two altitudes
$h_i^{\mathrm{baro}}$ and $h_i^{\mathrm{GNSS}}$. A record that does not follow the
fixed layout of the format, or whose time of day or position is impossible, is
discarded. A missing or unreadable altitude,
including the zero the format uses for a missing reading, is marked as missing, and
the fix keeps its time and position.

Time is measured in seconds from the first fix. Records carry only the time of day,
so a flight crossing UTC midnight would jump back by a day; a day is added when the
clock passes from the last hour of one day to the first hour of the next. Other
backward steps and repeated times are left in place, as are well-formed but implausible fixes such as position
spikes: the parser discards only what cannot be read. The validity flag, which marks
fixes without a three-dimensional GNSS solution, is kept but not used at this stage.

\section{Choice of altitude channel}
\label{sec:altchannel}

Some recorders log no pressure altitude, so we use GNSS altitude for every flight,
$z=h^{\mathrm{GNSS}}$. Its vertical datum is not uniform across recorders: the IGC
specification prescribes height above the WGS84 ellipsoid~\cite{fai_igc_spec}, the
CIVL specification for hang gliders and paragliders height above the
geoid~\cite{fai_civl_s7h_2024}. A constant offset cancels in altitude increments, but
comparing absolute altitudes across flights would need a datum audit.

Let $N$ be the number of fixes of a flight and $P$ the set of those with a GNSS
altitude. The flight is kept only if
\begin{equation}
  \frac{|P|}{N}\geq0.95
  \qquad\text{and}\qquad
  \max_{i\in P}h_i^{\mathrm{GNSS}}-\min_{i\in P}h_i^{\mathrm{GNSS}}\geq\SI{30}{\metre}.
  \label{eq:gnss-admission}
\end{equation}
The first condition rejects incomplete channels; the second rejects channels stuck at
one value, and flights with almost no vertical excursion. Fixes flagged as lacking a
three-dimensional solution count as present if they carry an altitude. Pressure
altitude is never substituted for a missing GNSS value.

The barometer is checked with the same two conditions and thresholds. A flight whose
barometer passes is marked, so that later cleaning steps can use pressure altitude as
independent evidence; the mark does not decide whether the flight is kept.
